%%
%% This is file `example_en.tex',
%% generated with the docstrip utility.
%%
%% The original source files were:
%%
%% ufrj-coppe.dtx  (with options: `exampleen')
%% 
%% This is one of the demonstration documents of the `ufrj' document class with
%% the `ufrj-coppe' unit style: the five main-language demos and the
%% side-by-side cover sheet.  It is generated from ufrj-coppe.dtx by docstrip;
%% edit the .dtx source, not this file.
%% 
%% Copyright (C) 2008-2026 Vicente Helano F. Batista, George O. Ainsworth Jr.,
%% Geraldo Xexéo and Eduardo Mangeli. Released under the GNU General Public
%% License version 3 (see the COPYING.txt file).
%% 
\documentclass[english,dsc]{ufrj}
\usepackage{ufrj-coppe}
\addbibresource{exemplo.bib}
\begin{document}
  \title{Título do Trabalho}
  \subtitle{um exemplo de uso do CoppeTeX em inglês}
  \foreigntitle{Title of the Work}
  \foreignsubtitle{an example of use of CoppeTeX in English}
  \author{Author's}{Name}
  \advisor{First}{Advisor}{D.Sc.}{UFRJ}
  \examiner{First Examiner}{D.Sc.}{UFRJ}
  \examiner{Second Examiner}{D.Sc.}{UFRJ}
  \department{PESC}
  \date{05}{2026}
  \keyword{Multilingual}
  \keyword{Demonstration}
  %% palavras-chave do resumo em português, a língua vernácula (3.1.2.1.4)
  \foreignkeyword{multilíngue}
  \foreignkeyword{demonstração}
  % Pre-textual order (UFRJ Manual 3.1.2), checked by the class: \maketitle,
  % \frontmatter, the abstract in Portuguese -- the vernacular language,
  % always the first one; here it is foreignabstract --, the abstract in
  % English (abstract), the lists and \tableofcontents, always the last.
  \maketitle
  \frontmatter
  \begin{foreignabstract}
    Este é o resumo em português, a língua vernácula, que vem antes do
    abstract qualquer que seja o idioma do trabalho (Manual 3.1.2). Num
    trabalho em inglês, quem o produz é o ambiente \texttt{foreignabstract}.
  \end{foreignabstract}
  \begin{abstract}
    This is the abstract in English, the main language of this work, used
    when the class is loaded with the \texttt{english} option. It comes after
    the abstract in Portuguese. Lorem ipsum dolor sit amet, consectetur
    adipiscing elit.
  \end{abstract}
  % Front-matter lists: only the table of contents is mandatory (UFRJ
  % Manual 3.1.2.1.6). The lists of figures, tables, abbreviations and
  % symbols are optional (3.1.2.2); the lists of quadros, programs and
  % algorithms appear when the work has one (COPPE norm, §10). This demo
  % has no illustrations, so it prints only the table of contents.
  \tableofcontents
  \mainmatter
  \chapter{Introduction}
  This chapter shows that the body is composed in English while the
  cover and folha-de-rosto remain in Portuguese (the UFRJ institutional
  template).
  \section{The deposit file: PDF/A}
  Deposit at UFRJ has been digital only since CEPG Resolution 246/2023, and the
  UFRJ Manual (2.2d) requires the final file in \textbf{PDF/A}, the long-term
  archiving format. PDF/A is PDF restricted by ISO~19005 for preservation: every
  font is embedded, colours declare a profile (sRGB) that makes them device
  independent, metadata go in an XMP packet that repositories read, and nothing
  depends on the outside --- no external content, JavaScript or encryption. The
  file then opens the same way decades from now, on any computer.

  The class uses \textbf{PDF/A-2b}: part 2 allows transparency, which part 1
  forbids and which coloured boxes and graphics often produce, and level
  \emph{b} guarantees faithful visual reproduction. The \texttt{pdfa} option in
  \verb|\documentclass[english,dsc,pdfa]{ufrj}| is all it takes: the class
  loads \texttt{pdfx}, builds the metadata from the title, author and keywords,
  and embeds the colour profile. Turn it on early, as \texttt{coppe-max-exemplo.tex}
  does: a non-conforming PDF or image shows up the day it is added, not the day
  before the deposit. To check, use the veraPDF validator, PDF/A-2b profile.

  \section{Write streams}
  pdf\LaTeX{} keeps at most sixteen files open for writing, and each list,
  index or glossary takes one. That is why the class loads
  \texttt{morewrites}. The \texttt{semmorewrites} option turns it off and saves
  about half a second per pass, but only suits a work with no index, no
  multiple indexes and no glossary.
\end{document}
%% 
%%
%% End of file `example_en.tex'.
